\subsection{GENERATING SERIES FOR THE APPELL POLYNOMIALS}

Let E be the ring of formal series in an indeterminate S, with coefficients in the ring of polynomials K{[}X{]} (Alg., IV. 41); in other words, the ring of formal series \(g(\mathrm{X}, \mathrm{S}) = \sum_{n=0}^{\infty} \alpha_n(\mathrm{X}) \mathrm{S}^n\) where the \(\alpha_n\) belong to K{[}X{]}. For every operator \(U\) on K{[}X{]} one defines a map \(U_{\mathrm{X}}\) of E to itself by putting \(U_{\mathrm{X}}(g(\mathrm{X}, \mathrm{S})) = \sum_{n=0}^{\infty} U(\alpha_n)\mathrm{S}^n\) . It is clear that E is a module over the ring K{[}{[}S{]}{]} of formal series in S with coefficients in K; by the linearity of \(U\) on K{[}X{]} one verifies immediately that for every element \(\theta \in \mathrm{K}[[\mathrm{S}]]\) and every \(g \in \mathrm{E}\) one has \(U_{\mathrm{X}}(\theta g) = \theta U_{\mathrm{X}}(g)\) ; in other words, \(U_{\mathrm{X}}\) is a linear map of the module E into itself.

PROPOSITION 6. Let \(U = D^{p}V = u(D)\) be a composition operator of order p on K{[}X{]}, \(u(S)\) being a formal power series of order p in K{[}{[}S{]}{]}. Then

\[
U _ {\mathrm{X}} \big (\exp (\mathrm{XS}) \big) = u (\mathrm{S}) \exp (\mathrm{XS})\tag{15}
\]

\[
\frac {\mathrm{S} ^ {p}}{u (\mathrm{S})} \exp (\mathrm{XS}) = \sum_ {n = 0} ^ {\infty} \frac {1}{n !} u _ {n} (\mathrm{X}) \mathrm{S} ^ {n}\tag{16}
\]

\(u_{n}\) being the Appell polynomial of index n attached to U.

By the scholium to th. 1 (VI, p.~271), to establish (15) it is enough to show that for every polynomial \(f(\mathrm{Y}) \in \mathrm{K}[\mathrm{Y}]\) one has

\[
U _ {\mathrm{X}} \left(\exp (\mathrm{XD} _ {\mathrm{Y}}) (f (\mathrm{Y}))\right) = u (\mathrm{D} _ {\mathrm{Y}}) \left(\exp (\mathrm{XD} _ {\mathrm{Y}}) (f (\mathrm{Y}))\right).\tag{17}
\]

Now the first term in (17) is \(U_{\mathrm{X}}(f(\mathrm{X} + \mathrm{Y}))\) , and, since \(U = u(\mathrm{D})\) , the second term in (17) is \(U_{\mathrm{Y}}(f(\mathrm{X} + \mathrm{Y}))\) , so the identity (17) reduces to (1) (VI, p.~270).

It now suffices to apply (15) to the composition operator \(V^{-1} = D^{p}/u(D)\) to obtain (16), since, by definition, one has

\[
V ^ {- 1} \big (\exp (\mathrm{XS}) \big) = \sum_ {n = 0} ^ {\infty} \frac {1}{n !} u _ {n} (\mathrm{X}) \mathrm{S} ^ {n}.
\]

Note that (16) can also be obtained by multiplying the formal series \(S^{p}/u(S)\) and \(\exp(XS)\) , taking account of (6).

One says that the formal series (16) is the generating series of the Appell polynomials attached to U.

